% !TeX root = ../main.tex
% Add the above to each chapter to make compiling the PDF easier in some editors.

\chapter{Literature Review}\label{chapter:literature-review}


\section{Patching}
Most vulnerabilities are not entirely novel; the knowledge required to address a new flaw is often already available in previously published fixes for similar issues. Since many software flaws recur across different programs or modules, existing patches can be automatically extracted, adapted, and applied to structurally similar code \cite{Shariffdeen2020}.
This strategy has been explored in Automated Program Repair, and prior work has shown how automated repair systems leverage existing defect and fix-pattern histories to generate patches for newly discovered bugs and vulnerabilities \cite{legoues2019}.

One early way to explore code to find vulnerabilities is to use tools like CodeQL and Joern. Such tools perform static analysis over a graph representation of the code, like encoding control flow, data flow, and syntactic structure to search for known patterns.
Even though these graph-based traversal approaches made vulnerability search much easier and more automated, they require expertise in the language and in the domain-specific vulnerabilities. To scale the search for vulnerabilities and avoid missing one because its pattern was not explicitly searched, it is necessary to shift from hand-crafted rules to an automated search.

\section*{Structural Retrieval and Graph Edit Distance}
Graph edit distance (GED) is defined as the minimum-cost sequence of vertex and edge insertions, deletions, and substitutions required to transform one graph into the other \cite{gao2010survey}.
This classical approach evaluates the structure directly and can be applied to arbitrary graph types, including as a measure of structural similarity between two graphs. This makes them a promising approach for identifying vulnerable structures already known in new code without any language as a proxy, and for reusing the graph knowledge directly in the search.
However, GED is NP-hard; additionally, exact/near-exact computation doesn't scale to corpus-level retrieval over thousands of historical fixes. 

This scalability issue is particularly relevant for security applications.
Feng et al. explore it in the context of vulnerability search, showing that control-flow-graph-based bug search does not scale to firmware corpora due to the cost of graph matching. Their approach converts CFGs into high-level numeric feature vectors that support real-time approximate
search via hashing \cite{feng2016scalable}. 
Xu et al. replace these hand-crafted features with a learned graph neural network embedding for
the same cross-platform bug-search task, reporting large accuracy gains over the hand-engineered baseline \cite{xu2017neural}. 
% the following work is interesting, but i didn't train any of the models so it can be used as a downside
% More generally, SimGNN reframes GED computation itself as a learning problem, training a neural network to approximate GED directly rather than computing it exactly \cite{li2019simgnn}. 
These works exemplify that structural similarity for security-relevant graphs is best operationalized by fixed-length embeddings whose distances approximate structural similarity, enabling standard approximate nearest-neighbor search over large historical corpora.
% (not by computing GED directly at query time, )


\section{Structural and Semantic Code Embedding Techniques}

Traditional graph classification frequently uses non-parametric graph kernels, offering computationally lightweight approaches that can run locally without requiring substantial overhead. 
One of its fields is structural comparison, whose core is the Graph Isomorphism problem, which determines whether the topologies of two graphs are identical. 
While no general polynomial-time algorithm for Graph Isomorphism is currently known \cite{garey1979computers, babai2016graph}, practical upper bounds on the computational complexity of the 1-dimensional Weisfeiler-Lehman (WL) color refinement test have been established \cite{weisfeiler1968reduction}.

The $1$-WL test characterizes graph topology through an iterative color-refinement process. 
At each iteration, a node aggregates its current label with the sorted multiset of its immediate neighbors' labels and hashes the resulting tuple into a new unique label. 
This process repeats until the label distribution between two compared graphs diverges or a fixed cutoff of $k$ iterations is reached; this latter only allows the isomorphism conclusion for the evaluated subgraph. 

The WL Fast Subtree Kernel \cite{shervashidze2011wl} applies this mechanism to a Reproducing Kernel Hilbert Space(RKHS), calculating the inner product of label histograms accumulated over $k$ iterations.
This approach establishes an expressivity benchmark for graph representations. However, despite its computational efficiency, $1$-WL is limited to local neighborhood tree structures defined by iterative local multiset hashing, preventing it from capturing higher-order global substructures such as cycles and cliques. 

The Graph Isomorphism Network (GIN) \cite{xu2019gin}, which bridges graph kernels and deep learning and is designed to maximize the discriminative power of Message-Passing Neural Networks, proved to be as powerful as the Weisfeiler-Lehman graph isomorphism tests. 
To achieve this performance, GIN uses learning neighborhood aggregation via a parameterized function, in which each node embedding iteratively encodes the structure of its k-hop neighborhood. 
To have the same power as 1-WL, GIN requires that both the neighborhood aggregation and graph-level readout functions must be strictly injective. 
A sufficient number of iterations is required to discriminate graphs; 
While deeper iterations are required for larger structural representations, features from early iterations may generalize better depending on the task.
% TODO: put the investigation of the number of iteration as future work

GraphFVD applies the graph embedding principles to the software security context.
It learns graph embeddings from the Code Property Graph code representation
Rather than embedding whole-function graphs, GraphFVD extracts property-graph slices (PrVCs) from the CPG and embeds them using a Hierarchical Attention Graph Convolutional Network.
Their experiments report that this slice-level representation outperforms both function-level and sequence-level slice baselines for vulnerability detection \cite{shao2025graphfvd}. 
This aligns with our design choice to embed a bounded slice rather than the full function CPG.
However, whereas GraphFVD's embeddings feed a supervised classifier for single-function detection, our approach indexes a historical corpus for a retrieval-augmented analysis.
Therefore, the two tasks share a representation, but their task objective diverge.

Complementing local message-passing architectures, the Network Laplacian Spectral Descriptor (NetLSD)\cite{netlsd} provides a global, multi-scale topological representation. 
NetLSD computes a heat-trace signature from the graph Laplacian's eigenvalues, yielding a permutation- and size-invariant embedding that requires no training and is comparable across graphs of different sizes \cite{netlsd}. It captures macro-level structural dynamics without requiring supervised training, making it computationally effective for comparing heterogeneous graph structures.

In contrast to purely structured embeddings, transformer-based models capture lexical and semantic regularities in source code. CodeBERT \cite{feng2020codebert} provides a pre-trained embedding that models token sequences, capturing variable naming conventions, API usage patterns, and control structure semantics—information that remains opaque to purely topological representations.

These embedding strategies of encoding code prioritize different aspects. WL and NetLSD are non-parametric structural descriptors that provide deterministic, invariant topological comparisons without training. GIN, just like GraphFVD, learns structural aggregation for specific vulnerability patterns. CodeBERT captures token-level semantics while lacking structural dependencies. 
Because no single representation captures both fine-grained topological flows and token-level semantics, combining structural graph embeddings with lexical representations offers a promising direction for vulnerability analysis and code retrieval.

\subsection*{Retrieval-Augmented Secure Code Generation}

RESCUE applies retrieval-augmented generation to secure code generation. It uses a hybrid knowledge base that combines vulnerability-fix data distilled into security guidelines with static program slicing to extract security-focused code examples. At retrieval time, RESCUE performs a hierarchical, multi-faceted lookup: first, it identifies the vulnerability type (CWE level) from the task description, then it retrieves fine-grained secure code examples within that type. 
RESCUE targets code generation for known vulnerability classes, while our pipeline targets the repair of already identified vulnerable functions. Both systems share key architectural insights by using program slicing to isolate security-relevant statements before embedding and employing hierarchical retrieval rather than flat nearest-neighbor lookup. Program slicing as a mechanism to concentrate security focus emerges as a convergent design pattern across recent work — alongside GraphFVD's property-graph-based vulnerability candidates and LLMxCPG's CPG-based slicing — reinforcing that isolating security-relevant subgraphs, rather than embedding whole functions, is a broadly validated approach for vulnerability-related tasks.

\section{Agentic Patching Frameworks}

LLM-based program repair has increasingly applied agentic frameworks that allow the model to iteratively gather context, interact with external tools, and revise its own output. (cite)

An example of this approach is AutoPatch, which leverages RAG-DB-assisted prompting to perform vulnerability detection and patching on a benchmark that mixes real-world and augmented Common Vulnerabilities and Exposures (CVEs).
The work evaluated the architecture against five popular LLMs, with GPT-4o achieving the best performance, yielding F1-scores of 90.3\% for vulnerability verification and 94.1\% for patching. The system performs particularly well in concurrency-related issues, validation, and logic handling. Furthermore, AutoPatch demonstrates greater cost efficiency than fine-tuning approaches with 10 epochs, incurring an approximately 1,209\% cost increase, while non-incremental fine-tuning at five CVE intervals resulted in a 5,230\% increase. The AutoPatch built benchmark serves as one of the two primary evaluation datasets utilized in this thesis (see Chapter~\ref{chapter:dataset}).
% To ensure the reliability of the generated patches, AutoPatch employs a rigorous patch-verification feedback loop. After a Code Patcher agent generates a candidate fix, the output is routed back to a Vulnerability Verifier agent. Using previously constructed one-shot examples, the Verifier analyzes the patched code for any remaining vulnerabilities. This cycle iterates until either no vulnerability is detected or a developer-specified maximum number of iterations is reached. While the average loop count was close to 1 across all models, certain difficult cases required more iterations (up to 8 for DeepSeek-R1 and 10 for o3-mini). These findings emphasize that context-aware verification—heavily reliant on in-context examples—is critical for capturing the semantic complexity of real-world vulnerabilities and surfacing logic inconsistencies.

Applying a similar tooling-based architecture, RepairAgent proposes an autonomous agent powered by Large Language Models to dynamically fix software bugs \cite{bouzenia2025repairagent}. 
With AutoPatch, RepairAgent represents a similar agentic generation proposed in this thesis in terms of task framing. 
The introduced agent plans and executes actions through a predefined set of rules that allows the LLM to gather information about the bug, gather repair ingredients, and validate fixes.
To evaluate these patches, the agent automatically applies the fix and executes the project's test suite, using the resulting test feedback to inform subsequent iterations. The agent decides to stop its iterative loop only when a patch successfully passes all validation tests or when a predefined computational budget is exhausted.
They evaluate their approach against Defects4J which contains 835 bugs .
% RepairAgent imposes an average cost of 270,000 tokens per bug, which, under the current pricing of OpenAI’s GPT-3.5 model, translates to 14 cents per bug. An additional evaluation on a set of recent bugs [25] shows that RepairAgent is able to achieve similar performance on single-line bugs while being a bit worse on multi-line and multi-file bugs, mainly, due to a higher complexity of bugs in GitBug-Java dataset.
% We believe from these results that RepairAgent is not much affected by potential data leakage of Defects4J.
% how does it decide when to stop, how does it evaluate patches.
% tool-augmented agentic loop, giving the LLM explicit actions for inspecting the codebase, retrieving relevant
% context, and validating candidate patches, and is the closest prior system to our own agentic generation layer in terms of task framing \cite{bouzenia2025repairagent}
% - ReAct — the general reasoning+acting loop RepairAgent (and your own agentic layer) build on. Cite as the pattern, not as a patching system itself.
% - Self-Refine — iterative self-critique loop; relevant if your agentic generation layer does any revision/retry step. If it doesn't, don't force this citation in — only include what you actually build on.
 
 QLCoder	agentic framework that automatically synthesizes CodeQL queries directly from CVE metadata. It bridges the gap between natural language vulnerability reports and actionable detection tools by embedding an LLM in an iterative synthesis loop that uses execution feedback to verify query correctness.	"RQ1: For how many real-world CVEs can QLCoder successfully generate syntactically valid and semantically precise queries?
 Structured Synthesis Loop: The agent proposes a query, a validator executes it on vulnerable and patched versions of a repository, and feedback is provided to refine the next proposal.
MFix Discrimination: Ensures the synthesized query is precise enough to flag the vulnerable version but remain silent on the patched version."	CWE-Bench-Java: 176 CVEs across 111 Java projects; 42 different CWE categories, such as Path Traversal (48), Cross-Site Scripting (36), and Code Injection (20); 65 CVEs reported in 2025 after Claude Sonnet 4 cutoff

MAAP is a self-evolving, role-specialized multi-agent framework designed for end-to-end automated vulnerability repair (AVR) in Python codebases. It addresses the limitations of single-agent approaches—such as context overload and myopic decision-making—by decomposing vulnerable repositories into fine-grained subtasks	
Contextual-Bandit Router: Learns to assign subtasks to the most appropriate specialized agents to optimize for success, cost, and latency.
% Retrieval-Augmented Knowledge Base: Stores successful and failed repair trajectories (AST diffs and rationales) to guide agents and prune ""dead-end"" explorations.
% Agent Factory: Dynamically spawns new specialized agents when the system encounters novel vulnerability patterns or low success rates"	VAITP selection: 120 real-world Python CVEs(Input validation and sanitization, memory corruption, authentication/session management). AI-generated and human-verified code snippets of vulnerable and patched versions
